\begin{textsample}{NEG Dim 4 – 2022\_c5 – Score -1.00 – 2022\_c5/enade\_2022\_tec\_proc\_ger\_587bf59e0a.txt}  \label{ex:f4_neg_077}
CAPÍTULO 6 DISTRIBUIÇÃO DE CONCEITOS A finalidade deste capítulo é apresentar um panorama nacional da distribuição do Conceito Enade 2022.
Para traçar esse panorama, recorre -se a tabelas e gráficos, que são acompanhados de comentários e análises que contribuem para a contextualização dos dados neles expostos.
Para visualizar o panorama do Enade/2002, em nível nacional, na primeira seção do capítulo, os dados sobre o número de cursos participantes são organizados por grande região.
Na segunda seção, o mapeamento do conceito do exame é feito por categoria administrativa das IES e modalidade de oferta dos cursos, de maneira cruzada às grandes regiões.
Na última seção, a lógica de apresentação dos dados é feita por organização acadêmica e, também, grande região de funcionamento dos cursos, no país como um todo. 6.1 PANORAMA NACIONAL DA DISTRIBUIÇÃO DOS CONCEITOS A lógica adotada na apresentação dos dados, nesta seção e nas duas seguintes, pressupõe o registro de duas informações de natureza reiterativa.
A primeira informação diz respeito aos valores possíveis do Enade, que variam em níveis de 1 a 5, sendo que à medida que esses valores aumentam, o curso recebeu uma avaliação melhor além disso, os cursos com somente um ou sem nenhum concluinte participante não obtêm o Conceito Enade, ficando Sem Conceito ( SC ), como já ressaltado na Apresentação neste relatório.
O segundo registro é que a diferença entre os cursos tabulados neste capítulo e no Capítulo 3 corresponde aos cursos que não foram avaliados, ou seja, àqueles, em princípio, sem estudantes concluintes que tivessem sido inscritos para a prova.
Na Tabela 6.1, são apresentadas a quantidade e a distribuição de cursos de Tecnologia em Processos Gerenciais participantes do Enade/2022, por faixa de conceito e grande região.
Observando -se os dados da Tabela 6.1, nota -se que, dos 327 cursos participantes, 124 ( 37,9 \% ) classificaram -se com conceito 3, o valor modal.
Esse foi também o conceito modal na maioria das regiões : Nordeste ( 26,5 \% ), Sudeste ( 34,5 \% ), Sul ( 45,2 \% ) e Centro-Oeste ( 47,1 \% ).
O conceito 4 foi o segundo mais frequente em nível nacional ( 20,2 \%, correspondendo a 66 cursos ), tendo sido o conceito modal na região Sudeste ( 23 \% ).
O conceito 2 foi o terceiro ( 16,2 \%, correspondendo a 53 cursos ).
Houve, ainda, 19 cursos ( 5,8 \% ) que receberam o conceito 5, e 5 cursos ( 1,5 \% ) que receberam o conceito 1.
Dos 327 cursos de Tecnologia em Processos Gerenciais, 60 ( 18,3 \% ) ficaram Sem Conceito ( SC ).
A região Norte participou com 7 cursos, o que corresponde a 2,1 \% do total nacional.
Desses, 2 cursos receberam o conceito 2, o que equivale a 28,6 \% do total regional.
O conceito modal na região foi 3, equivalente a 42,9 \% da região. 1 curso ( 14,3 \% ) atribuiu -se o conceito 1.
A região Nordeste participou com 34 cursos, o correspondente a 10,3 \% do total nacional.
Nessa região, 9 cursos ( 26,5 \% em termos regionais ) obtiveram conceito 3, outros 9 cursos receberam conceito 2 ( 26,5 \% ), configurando assim os conceitos 2 e 3 em categorias modais.
O conceito 1 foi atribuído a 1 curso ( 2,9 \% ).
Ainda nessa região, 4 cursos ( 11,8 \% ) foram avaliados com o conceito 4, e 3 cursos ( 8,8 \% ) foram avaliados com o conceito 5.
Oito cursos ficaram Sem Conceito ( SC ).
Para os 165 cursos participantes da região Sudeste ( 50,4 \% do total nacional ), o conceito 3 foi o conceito modal, como já anteriormente assinalado, tendo sido obtido por 57 cursos ( 34,5 \% ).
O conceito 1 foi atribuído a 3 cursos ( 1,8 \% ), e o conceito 2 também foi atribuído, a 27 cursos ( 16,3 \% ).
Receberam os conceitos 4 e 5, respectivamente, 38 cursos ( 23 \% ) e 10 cursos ( 6,1 \% ).
Nessa região, 30 cursos ficaram Sem Conceito ( SC ).
Os 104 cursos da região Sul corresponderam a 31,8 \% do total nacional.
Como já indicado, houve predominância do conceito 3 ( conceito modal ), atribuído a 45,2 \% dos cursos da região, ou seja, atribuído a 47 dos 104 cursos participantes na região Sul.
O conceito 1 foi atribuído a um curso ( 1 \% ) ; o conceito 2, a 14 cursos ( 13,5 \% ) ; e os conceitos 4 e 5, a 23 ( 22,1 \% ) e a 3 cursos ( 2,9 \% ), respectivamente.
Na região Sul, 16 cursos ficaram Sem Conceito ( SC ).
Já dos 17 cursos participantes na região Centro-Oeste ( 5,1 \% do total nacional ), 8 ( 47,1 \% em termos regionais ) receberam conceito 3, o conceito modal, como já destacado. 2 cursos ( 11,7 \% ) obtiveram o conceito 5, um curso ( 5,8 \% ) obteve o conceito 4, e o mesmo aconteceu com o conceito 2, somente 1 curso.
Nessa região, 5 cursos ficaram Sem Conceito ( SC ).
Em um gráfico de distribuição cumulativa, em cada valor de abscissa ( eixo x ), o valor apresentado no eixo das ordenadas ( eixo y ) é o valor acumulado da variável ( Conceito Enade ).
Em linhas gerais, em um gráfico desta mesma tipologia, poligonais mais à esquerda representam, no caso dos dados apresentados nesta seção, as grandes regiões com a pior distribuição, e poligonais mais à direita, as grandes regiões com os melhores conceitos.
A partir destas diretrizes, no Gráfico 6.1, é apresentada a distribuição cumulativa do Conceito Enade, segundo a grande região.
Nesse gráfico, poligonais mais à esquerda representam valores mais baixos do Conceito Enade, e poligonais mais à direita, valores mais altos. 6.2 CONCEITOS POR CATEGORIA ADMINISTRATIVA, MODALIDADE DE OFERTA DOS CURSOS E GRANDE REGIÃO Na Tabela 6.2, é apresentada a distribuição dos cursos de Tecnologia em Processos Gerenciais participantes do Enade/2022, por categoria administrativa e modalidade de oferta, de acordo com os conceitos por eles alcançados, segundo as grandes regiões brasileiras.
De acordo com as informações da tabela 6.2, dos 327 cursos participantes, 59 ( 18 \% ) eram ministrados em instituições públicas, e 268 ( 81,9 \% ), em instituições privadas.
Quanto à modalidade de oferta, 159 cursos eram a distância, e 168 eram, presenciais.
Ainda conforme as informações da mesma tabela, em termos nacionais, entre instituições públicas e privadas observa -se que, dos 19 cursos avaliados com conceito 5, 9 foram oferecidos em IES públicas, e 10, em IES privadas.
Dos 59 cursos participantes de IES públicas, o conceito 4 foi o valor modal, atribuído a 30 cursos ( 50,8 \% ).
Dos demais cursos de instituições públicas participantes, nenhum dos cursos foram avaliados com o conceito 1, 3 cursos ( 5 \% ) foram avaliados com o conceito 2, 17 cursos ( 28,8 \% ) foram avaliados com o conceito 3, e 9 cursos ( 15,2 \% ), com o conceito 5.
Nessa categoria, nenhum curso ficou Sem Conceito ( SC ).
Na rede privada, constituíram conceito modal foi o 3, atribuído a 107 cursos, o correspondente a 39,9 \% dos 268 cursos da categoria.
Dos demais cursos participantes, 5 ( 4,3 \% ) receberam o conceito 1, e 50 ( 18,6 \% ), o conceito 2.
O conceito 4 foi atribuído a 36 cursos ( 13,4 \% ), e o conceito 5, a 10 cursos ( 3,7 \% ).
Nessa categoria administrativa, 60 cursos ( 22,3 \% ) ficaram Sem Conceito ( SC ).
Dos 168 cursos participantes da modalidade presencial, o conceito 4 foi o valor modal, atribuído a 53 cursos ( 31,5 \% ).
Dos demais cursos nessa modalidade de oferta, 47 cursos ( 27,9 \% ) receberam o conceito 3 ; 22 cursos ( 13 \% ) recebeu o conceito 2 ; 9 cursos ( 5,3 \% ) receberam conceito 5 ; e apenas 1 curso recebeu conceito 1.
Nessa modalidade de oferta, 36 cursos ( 21,4 \% ) ficaram Sem Conceito ( SC ).
Dos 159 cursos participantes da modalidade a distância, o conceito 3 foi o valor modal, atribuído a 77 cursos ( 48,4 \% ).
Dos demais cursos nessa modalidade de oferta, 4 cursos ( 6,8 \% ) receberam o conceito 1 ; 31 cursos ( 36,8 \% ) recebeu o conceito 2 ; 13 cursos, ( 8,1 \% ), o conceito 4 ; e 10 cursos ( 6,2 \% ), o conceito 5.
Nessa modalidade de oferta, 24 cursos ( 15 \% ) ficaram Sem Conceito ( SC ).
Na análise por região, observa -se que, na região Norte, as instituições privadas participaram com 7 cursos ( 2,1 \% do total regional ), e 1 curso ficou Sem Conceito ( SC ).
O conceito modal para as instituições privadas na região foi o 3, atribuído a 3 cursos, correspondendo a 42,8 \% dessa categoria na região.
Nenhum curso foi avaliado com o conceito 1 ; 1, com o conceito 5 ; e nenhum, com conceito 4.
As instituições públicas não participaram na região Norte.
Na região Nordeste, a rede privada concentrou 31 dos 34 cursos participantes, o equivalente a 91,1 \% do total da região.
Entre os cursos oferecidos por IES privadas no Nordeste, 9 cursos foram avaliados com conceito 3, o conceito modal. 1 curso obteve o conceito 1 ; 8, o conceito 2 ; 3, o conceito 4 ; e 2 cursos obtiveram o conceito 5.
Nessa categoria administrativa e grande região, 8 cursos ficaram sem Conceito ( SC ).
As instituições públicas da região Nordeste participaram com 3 cursos ( 8,82 \% do total da região ).
Desses, 1 ( 33,3 \% ) obteve conceito 2 ; seguido por 1 curso, que recebeu o conceito 4 ; e 1, o conceito 5.
Nenhum curso ficou Sem Conceito ( SC ).
Dos cursos oferecidos na região Nordeste, 15 eram presenciais : 5 ficaram sem conceito ( SC ) ; 3, com o conceito 2 ; 2, com o conceito 3 ; 3, com o conceito 4 ; e 2, com o conceito 5.
Na modalidade a distância, eram 19 cursos, tendo um obtido o conceito 1 ; 6, o conceito 2 ; 7, o conceito 3, o valor modal ; e 1, o conceito 4.
Nessa região, 3 cursos ficaram Sem Conceito ( SC ).
Na região Sudeste, a proporção de cursos da rede privada, 75,5 \%, correspondeu a 125 dos 165 cursos participantes.
Nessa combinação de categoria administrativa e grande região, o conceito modal foi o 3, atribuído a 57 cursos. 30 cursos ficaram Sem Conceito ( SC ).
Os demais foram avaliados com os seguintes conceitos : 1 ( 3 cursos ), 2 ( 27 cursos ), 4 ( 38 cursos ) e 5 ( 10 cursos ).
Dos 40 cursos oferecidos por instituições públicas na região Sudeste, o conceito modal foi o 4, atribuído a 21 cursos.
Dos demais cursos, 12 receberam o conceito 3 e 6 receberam o conceito 5.
Também na região Sudeste, grande parcela dos cursos foi oferecida na modalidade presencial : 90 de 165 cursos.
Para esses, a moda foi o conceito 4, com 30 cursos.
Os demais cursos presenciais da região foram \textbf{alocados} nos conceitos, 2 ( 11 cursos ), 3 ( 26 cursos ) e 5 ( 5 cursos ).
Ficaram Sem Conceito ( SC ) 18 cursos presenciais na região Sudeste.
Na região Sudeste, 75 cursos foram oferecidos na modalidade a distância, tendo sido 3 cursos com conceito 1 ; 16 com conceito 2 ; 31 com conceito 3, o valor modal ; 8 com o conceito 4 ; e 5 com o conceito 5.
Na região Sudeste, 12 cursos a distância ficaram Sem Conceito ( SC ).
As instituições privadas concentraram 268 dos 91 cursos participantes da região Sul, 33,9 \% do total regional.
Desses, 42 cursos obtiveram conceito 3, o conceito modal.
Nessa combinação de categoria administrativa e grande região, 1 curso recebeu o conceito 1 ; 14 receberam o conceito 2 ; outros 16, o conceito 4 ; e ainda 2, o conceito 5. 16 cursos ficaram sem Conceito ( SC ).
As instituições públicas na região Sul participaram com 13 cursos ( 4,8 \% dos cursos da região ), aos quais foram atribuídos os conceitos 5 ( 1 curso ), 3 ( 5 cursos ) e 4 ( 7 cursos ), sendo este último o modal do grupo.
Nenhum curso da região Sul oferecido por IES públicas ficou Sem Conceito ( SC ).
Na região Sul, 55 dos cursos foram oferecidos na modalidade presencial.
Para esses, a moda foi o conceito 4, com 19 cursos.
Os demais cursos presenciais da região foram \textbf{alocados} nos conceitos 1 ( um curso ), 2 ( 7 cursos ), 3 ( 17 cursos ) e 5 ( um curso ).
Ficaram sem Conceito ( SC ) dez cursos.
Na modalidade a distância, foram oferecidos 49 cursos, sendo 7 com conceito 2, 30 com conceito 3, o valor modal, 4 com o conceito 4.
Seis cursos ficaram sem Conceito ( SC ).
Na região Centro-Oeste, 14 dos 17 cursos participantes eram de IES privadas ( 82,3 \% em termos regionais ).
Desses, 8 cursos receberam conceito 3, o valor modal.
Dos demais cursos, um recebeu o conceito 5.
Cinco cursos ficaram Sem Conceito ( SC ).
Quanto aos 3 cursos oferecidos por instituições públicas na região Centro-Oeste ( 17,5 \% do total regional ), 1 apresentou o conceito 3 ; outro, o 4 ; e o último, 5.
Nenhum curso ficou Sem Conceito ( SC ) nessa região.
Também na região Centro-Oeste, a maioria dos cursos foi a distância : 10 ( 58,2 \% ).
Para esses, a moda foi o conceito 3, com 7 cursos ( 70 \% ).
O outro curso a distância da região foi \textbf{alocado} no conceito 5 ( um curso, 14,9 \% ). 3 cursos presenciais da região ficaram sem Conceito ( SC ).
A título de síntese, no Gráfico 6.2, apresenta -se a distribuição do Conceito Enade, segundo a categoria administrativa da IES.
Nesse gráfico, como dito anteriormente, poligonais mais à esquerda representam valores mais baixos do Conceito Enade, e poligonais mais à direita, valores mais altos.
É importante observar que, para cada poligonal presente nesse gráfico, a porcentagem constante no eixo y se refere ao universo de cursos oferecidos por instituições públicas ( poligonal vermelha ) ou privadas ( poligonal azul ). 6.3 CONCEITOS POR ORGANIZAÇÃO ACADÊMICA E GRANDE REGIÃO Na Tabela 6.3, encontra -se a distribuição dos conceitos atribuídos aos cursos participantes do Enade/2022, na área de Tecnologia em Processos Gerenciais, por organização acadêmica, segundo as grandes regiões brasileiras.
Como mostram os dados da Tabela 6.3, dos 327 cursos de Tecnologia em Processos Gerenciais participantes, 89 eram oferecidos em universidades ; 119, em centros universitários ; 97, em faculdades ; e 22, em CEFET/IF.
Esta distribuição corresponde a, respectivamente, 27,2 \%, 36,3 \%, 29,6 \% e 6,7 \% dos cursos.
De acordo com os dados apresentados na mesma tabela, de todos os 19 cursos avaliados com o conceito 5, 7 eram vinculados a universidades.
Os cursos vinculados a esse tipo de organização acadêmica tiveram o conceito 3 como conceito modal, atribuído a 124 cursos ( 37,9 \% ).
Os demais cursos ofertados por universidades receberam os conceitos 1 ( nenhum curso ), 2 ( 13 cursos ), 4 ( 17 cursos ) e conceito 5 ( 7 cursos, como já mencionado ). 15 cursos ficaram Sem Conceito ( SC ).
Dos cursos vinculados a centros universitários, o conceito modal também foi o 3, atribuído a 53 cursos ( 44,5 \% ). 26 cursos vinculados a esse tipo de organização acadêmica ficaram Sem Conceito ( SC ), e os outros cursos receberam os conceitos 1 ( 2 cursos ), 2 ( 23 cursos ), 4 ( 10 cursos ) e 5 ( 5 cursos ).
Dos 97 cursos mantidos por faculdades, 27 ( 27,8 \% ) receberam conceito 3, o conceito modal.
O segundo conceito mais frequentemente atribuído aos cursos oferecidos por essa categoria de IES foi o conceito 4, recebido também por 27 cursos.
Na sequência, 15 cursos obtiveram conceito 2 ; 6 cursos, o conceito 5.
Dos cursos oferecidos em faculdades, 19 ficaram Sem Conceito ( SC ).
Dos 22 cursos mantidos por CEFET/IF, doze ( 54,5 \% ) receberam conceito 4, o conceito modal.
O segundo conceito mais frequentemente atribuído aos cursos oferecidos por essa categoria de IES foi o conceito 3, recebido por sete cursos.
Na sequência, dois cursos obtiveram o conceito 2, nenhum curso ficou Sem Conceito ( SC ) e um curso, com o conceito 5.
Considerando -se, separadamente, as regiões brasileiras, verifica -se que, na região Norte, as universidades participaram com 3 cursos avaliados, aos quais foram atribuídos os seguintes conceitos : 2, a 1 curso ; 3, a 1 curso ; 5, a 1 curso.
Não foi atribuído curso Sem Conceito ( SC ) nessa combinação de categoria e região.
Os centros universitários da região Norte foram representados por três cursos, dos quais um obteve conceito 2 e um conceito 3.
Os conceitos 1, 4 e 5 não foram atribuídos a nenhum curso.
Do mesmo modo, a um curso foi atribuído Sem Conceito ( SC ).
As faculdades participaram com 1 curso na região Norte, o qual recebeu o conceito 3 ( um curso ).
Nenhum curso ficou Sem Conceito ( SC ).
Nessa região e organização acadêmica, nenhum curso foi oferecido em CEFET/IF.
Na região Nordeste, as universidades participaram com 34 dos 327 cursos da Área de Tecnologia em Processos Gerenciais oferecidos.
O conceito modal foi o 2 e o 3, atribuídos a 9 cursos cada. 4 cursos obtiveram o conceito 4 ; três cursos, o conceito 5.
Um curso oferecido por universidades no Nordeste recebeu o conceito 1 e 8 cursos ficaram Sem Conceito ( SC ).
Os centros universitários contaram com 20 cursos participantes na região Nordeste, 12 dos quais receberam o conceito 2 e 3 modal, sendo 6 cursos para cada.
Para os demais cursos, foram atribuídos os conceitos 4 ( um curso ) e 5 ( um curso ).
Seis cursos oferecidos por centros universitários no Nordeste ficaram Sem Conceito ( SC ).
As faculdades foram representadas por 6 cursos na região Nordeste, desses 6 cursos, tivemos um curso para os conceitos 1, 2, 3 e 4.
Dois cursos ficaram sem Conceito ( SC ).
Na organização acadêmica CEFET/IF, tivemos a participação de 3 cursos : 1 curso para o conceito 2 1 curso para o conceito 4 e outro curso para o conceito 5.
Na região Sudeste, as universidades concentram 45 dos 165 cursos de Tecnologia em Processos Gerenciais da região.
Dos cursos oferecidos em IES com este tipo de organização acadêmica, nessa região, o conceito modal foi 3, atribuído a 15 cursos.
Os demais cursos receberam os conceitos 2 ( 8 cursos ), 4 ( 10 cursos ), 5 ( 2 cursos ) e 2 ( 8 cursos ). 30 cursos ficaram Sem Conceito ( SC ), e nenhum recebeu o conceito 1.
Os centros universitários participaram com 61 cursos na região Sudeste, dos quais 27 obtiveram o conceito modal 3.
O conceito 1 foi atribuído a dois cursos ; o conceito 2, a 12 cursos ; o conceito 4, a 6 cursos ; e o conceito 5, a dois cursos.
Doze cursos ficaram Sem Conceito ( SC ) nessa combinação de organização acadêmica e região.
Já as faculdades foram representadas por 52 cursos na região Sudeste, dos quais oito ficaram Sem Conceito ( SC ).
O conceito modal 4 foi atribuído a 18 cursos.
Aos demais cursos foram atribuídos os conceitos 1 ( um curso ), 2 ( 7 cursos ), 3 ( 12 cursos ) e 5 ( seis cursos ).
Para os sete cursos participantes de CEFET/IF, na região Sudeste, os conceitos foram : 3 ( três cursos ), 4 ( quatro cursos ), o valor modal.
Nenhum curso ficou Sem Conceito ( SC ) ou recebeu os conceitos 1, 2 e 5 nesse tipo de organização acadêmica.
Dos 104 cursos da região Sul, 31 eram mantidos por universidades, para os quais o conceito modal foi 3, atribuído a 16 cursos.
Os demais cursos receberam os conceitos 2 ( 3 cursos ), 4 ( 6 cursos ), 5 ( um curso ) ; nenhum curso recebeu o conceito 1.
Em IES com esse tipo de organização acadêmica, cinco cursos da região Sul ficaram Sem Conceito ( SC ).
Dos 29 cursos participantes de centros universitários da região Sul, quatro ficaram Sem Conceito ( SC ), e 16 cursos receberam conceito 3, o valor modal.
Os conceitos restantes ficaram assim distribuídos : nenhum curso ( conceito 1 ) ; 4 cursos ( conceito 2 ) ; três cursos ( conceito 4 ) ; e dois cursos ( conceito 5 ).
Quanto aos 34 cursos vinculados a faculdades na região Sul, 11 cursos receberam o conceito 3, valor modal, e sete cursos ficaram Sem Conceito ( SC ).
Os demais cursos receberam os conceitos 1 ( um curso ), 2 ( 7 cursos ), 4 ( oito cursos ).
Dos 10 cursos oferecidos em CEFET/IF, seis cursos foram avaliados com o conceito 4, sendo modal.
Outros 4 cursos receberam o conceito 3.
Nesse tipo de organização, nenhum curso da região Sul ficou Sem Conceito ( SC ) ou com os conceitos 1, 2 e 5.
Na região Centro-Oeste, 5 dos 17 cursos oferecidos eram mantidos por universidades.
Para esse tipo de organização acadêmica, o conceito modal foi 3, atribuído a 3 cursos ; e, nenhum curso da região ficou Sem Conceito ( SC ).
Os demais cursos ( 2 ) receberam os conceitos 5.
Os centros universitários da região Centro-Oeste contaram com 6 cursos : 3 cursos foram avaliados com o conceito 3 ( modal ) ; e três cursos ficaram Sem Conceito ( SC ).
Dos 4 cursos oferecidos em faculdades na região Centro-Oeste, 2 cursos receberam o conceito 3 ( modal ) ; e dois cursos oferecidos em faculdades ficaram Sem Conceito ( SC ).
Dos dois cursos oferecidos em CEFET/IF, dois cursos foram avaliados, um com o conceito 2 e outro com o conceito 4.
Nesse tipo de organização acadêmica, nenhum curso da região Centro-Oeste ficou Sem Conceito ( SC ).
De maneira semelhante à forma como se procedeu em relação à seção anterior, encerra -se com o Gráfico 6.3, que apresenta a distribuição do Conceito Enade dos cursos de Tecnologia em Processos Gerenciais segundo a organização acadêmica da IES.
A leitura desse último gráfico é semelhante à leitura feita dos anteriores, seguindo -se a lógica de que poligonais deslocadas para a esquerda correspondem a valores mais baixos do Conceito Enade ( 1, 2 e 3 ) e poligonais deslocadas para a direita correspondem a valores mais altos do Conceito Enade.

% matched lemmas: alocar
\end{textsample}
